\section{Removing the fixed-point-free involution loss in the alternating gate}
\label{sec:ordinary-alternating-fixed-free}

\paragraph{Status and scope.}
This complete new root derivation has passed the full root and
independent /root/g34\_eightfolds proof reviews, with separate exact
arithmetic supplements. Its new published input is Fomin--Lulov's exact
rim-hook estimate. The paper's actual theorem and complete proof in
Sections1--2, pp.~219--223, have been read; no reconstruction of its
upstream combinatorial inputs is claimed. The arguments below concern
the same transitive ordinary alternating core and original imaginary
character-field inclusion as HD-ALT-04--06.

\begin{lemma}[HD-ALT-07a: the three-cycle bound suffices for all long shapes]
\label{HD-ALT-07a}
Let $n\ge15$, let $\lambda=(n-k,\mu)$ with $\mu\vdash k$ and
$1\le k<n/4$, and let $W=S^\lambda|_{A_n}$ have degree $a$.
Every nonscalar finite-order linear normalizer $T$ of this represented
core satisfies
\[
 \nu(T)\ge a\beta(n,k),\qquad
 \beta(n,k)=\frac{k(n-1-k)}{(n-1)(n-2)}<\frac14.
 \tag{ALTFF.1}
\]
All scalar phases and all compatible outer automorphisms are allowed.
\end{lemma}

\begin{proof}
The long first row prevents $\lambda$ from being self-conjugate, so its
alternating restriction is irreducible. Every normalizer is scalar
times the operator of an element $g\in S_n$, by the full-automorphism
and Schur-lemma argument of HD-ALT-03a. A nonscalar normalizer has
$g\ne1$. Multiplication by its scalar phase preserves $\nu$.

If $k=1$, this is the deleted standard module. Every nonidentity
permutation has a nonscalar operator on it and hence support at least
one, which equals $a\beta(n,1)$ since $a=n-1$.
Now assume $k\ge2$. The precise degree consequence used in
HD-ALT-04--06 gives
\[
                         a\ge\frac{n(n-3)}2>4(n+1).
 \tag{ALTFF.2}
\]
Here the two tail-two degrees are $n(n-3)/2$ and
$(n-1)(n-2)/2$; every other nonstandard long shape has
$a\ge2n^3/27\ge n(n-3)/2$. These are the same published
low-degree classification consequences, with locator
\cite[proof of Theorem~9.7, p.~312]{HD:GKT2019}.
For the strict inequality in (ALTFF.2), twice the difference is
$n^2-11n-8$; at $n=15+x$ this is $x^2+19x+52>0$.

The commutator construction of HD-ALT-03a produces a three-cycle
unless $g$ is a fixed-point-free involution. In every case with that
three-cycle, the exact rank inequality
$\operatorname{rank}([A,B]-I)\le2\nu(A)$ and equality of the two
nontrivial three-cycle eigenmultiplicities give
\[
 \nu(T)\ge\frac{1-u}{3}a\ge\beta(n,k)a,
 \qquad u=\chi_W((123))/a,
 \tag{ALTFF.3}
\]
where the last inequality is the proved tail-content estimate
HD-ALT-05a. The transposition in the commutator acts in the actual symmetric
extension; no compatible outer action is assumed on a split constituent.

It remains that $n=2m$ and $g$ has cycle type $2^m$.
Fomin--Lulov's published Theorem1.1, with the character interpretation
immediately following it on p.~220, states
\[
 |\chi^{\lambda}(g)|\le
       \frac{m!\,2^m}{\sqrt{n!}}\sqrt a.
 \tag{ALTFF.4}
\]
See \cite[Theorem~1.1 and following character equality, p.~220;
complete proof, pp.~221--223]{HD:FominLulov1995}.
If the shape is not domino-decomposable, its value is zero by the
same rim-hook rule, so the displayed bound still holds.
The central binomial coefficient is the largest coefficient of
$(1+x)^{2m}$, hence
\[
 \binom{2m}{m}\ge\frac{4^m}{2m+1},\qquad
 \left(\frac{m!2^m}{\sqrt{n!}}\right)^2
     =\frac{4^m}{\binom{2m}{m}}\le n+1.
 \tag{ALTFF.5}
\]
Consequently (ALTFF.2) gives
\[
 \frac{|\chi^\lambda(g)|}{a}
      \le\sqrt{\frac{n+1}{a}}<\frac12,
 \qquad
 \nu(g|W)=\frac{a-|\chi^\lambda(g)|}{2}>\frac a4.
 \tag{ALTFF.6}
\]
The last equality uses the two eigenvalues of an involution and
continues to compute projective support after any scalar phase.
It applies whether $g$ is even or odd, since $W$ has its stated
symmetric extension and the underlying vector space is unchanged.

Finally $k(n-1-k)$ is increasing for $1\le k<n/4$, and
\[
 \beta(n,k)<\frac{n(3n-4)}{16(n-1)(n-2)}<\frac14.
 \tag{ALTFF.7}
\]
The final inequality is equivalent to $n^2-8n+8>0$, whose value
at $15+x$ is $x^2+22x+113$. Thus (ALTFF.6) is stronger than
(ALTFF.1) and covers the only remaining case.
\end{proof}

\begin{theorem}[HD-ALT-07b: an exact longer-row threshold]
\label{HD-ALT-07b}
In the transitive ordinary alternating-core setting of HD-ALT-04b,
for every remaining $n\ge15$ factor label one has
\[
 \max(\lambda_1,\lambda'_1)>
                \frac{1+\sqrt{3/5}}2\,n.
 \tag{ALTFF.8}
\]
More precisely, on writing the long-row label as $(n-k,\mu)$ and
letting $e$ be the actual total inducing degree, the necessary
finite-degree inequality is
\[
 \frac{k(n-1-k)}{(n-1)(n-2)}
                   \le\frac1{10}+\frac3{5e}.
 \tag{ALTFF.9}
\]
The assertion includes compatible outers, scalar phases, central scalar
extensions and all factor permutations. It does not assert the existence
of any source satisfying these inequalities.
\end{theorem}

\begin{proof}
By the already proved three-quarter bound, transpose the label if
necessary and write $\lambda=(n-k,\mu)$ with $0<k<n/4$.
Transitivity gives the same label, up to transpose, on every factor.
The factor representation is real, so the tensor core has a unique
nondegenerate symmetric-form line. Its multiplier on the actual
inducing group is nontrivial: otherwise the character would be real,
contrary to the inclusion of the original imaginary quadratic field.
The actual intermediate-cover and period-factorization arguments give
genus zero, at least twelve genuine inertias, and total moving rank
$2e+6$, as in HD-ALT-06b. At least two inertias have nontrivial form
multiplier, each moving at least $e/2$ dimensions.

The support fraction (ALTFF.1) passes to factorwise tensor operators
by spectral convolution. Every nontrivial factor-permutation operator
has the stronger quarter bound, since $\beta(n,k)<1/4$.
A genuine nonidentity scalar moves the entire space. Applying these
bounds to two multiplier inertias and ten others gives
\[
                         2e+6\ge e+10e\beta(n,k).
 \tag{ALTFF.10}
\]
This proves (ALTFF.9), with the actual inducing field and actual
multiplicity variation throughout.

Put
\[
 c=\frac{1-\sqrt{3/5}}2,
 \qquad c(1-c)=\frac1{10},\qquad \frac1{10}<c<\frac18.
 \tag{ALTFF.11}
\]
Suppose (ALTFF.8) fails, so $k\ge cn$.
Since $cn>1$ for $n\ge15$, the integer $k$ is at least two and
(ALTFF.2) applies. Monotonicity in $k$ gives
\[
 \begin{split}
 \beta(n,k)-\frac1{10}
 &\ge\beta(n,cn)-c(1-c)\\
 &=\frac{c\{(2-3c)n-2+2c\}}{(n-1)(n-2)}
   >\frac{13n-16}{80n^2}.
 \end{split}
 \tag{ALTFF.12}
\]
Indeed $c>1/10$, $2-3c>13/8$, $-2+2c>-2$, and the positive
denominator is less than $n^2$. The lower bracket $13n/8-2$ and the factor $c$ are positive
for $n\ge15$, which justifies multiplying these strict inequalities.
Furthermore
\[
 \frac{13n-16}{80n^2}
       >\frac6{5n(n-3)}\ge\frac3{5e}.
 \tag{ALTFF.13}
\]
The first inequality is equivalent to
$13n^2-151n+48>0$; at $n=15+x$ this polynomial is
$13x^2+239x+708>0$. The second follows from
$e\ge a\ge n(n-3)/2$. Equations (ALTFF.12)--(ALTFF.13)
contradict (ALTFF.9). This proves the exact threshold for all
$n\ge15$, with no asymptotic exception or finite table extrapolation.
\end{proof}

\paragraph{Strictness and framework.}
For example $(88,12)$ at $n=100$ has longest row $88>87n/100$
but $88<\tfrac12(1+\sqrt{3/5})n$. Moreover the stronger previous
finite-degree gate is also satisfied: HD-ALT-06a gives
$\delta(100,12)=1519571/15684900<1/10$, whereas the new bound is
$\beta(100,12)=58/539>1/10$. Thus this is a strict reduction even
relative to the previous explicit inequality, not only its uniform
$87/100$ consequence. Its two-row degree is
$\binom{100}{12}-\binom{100}{11}=908791246463100$.
No existence of an actual source with that label is asserted.
This is a restriction on actual finite local systems and weight-one
period maps, hence on condition-A suppliers (cote16 scan050;
SIGN\_NEXT SN-C03/C04). It does not produce cycles through condition I
(scan049 and cote44 Theorem4), change the resultant/Cayley--Hamilton
projector of scans048--049, or add an $L$-value or weight premise.
Its scheme, representation-ring, lambda/K-theory, EGA/SGA/FGA and
motivic placements are the common framework already stated.
The still longer shapes, including any standard tensor cases not
separately excluded, remain open. Markman, Severn and032 keep their
previous conditional status wherever used.
